\thispagestyle{plain}
\noindent{\sffamily\heiti\Large KSSOLV Toolbox 中文使用者手冊}

\vspace{5mm}
\noindent 獨立應用程式、MATLAB\regmark{} 工具箱與\ksonline{}\par
\noindent 適用版本：V0.5.0\quad 2026 年 10 月

\vspace{7mm}
本書介紹 KSSOLV Toolbox 的安裝、結構準備、電子結構計算和結果分析。教程從小型分子與晶體出發，逐步延伸到表面、異質結、扭轉結構、低維模型與軟物質。本書適用於獨立應用程式、MATLAB\regmark{} 工具箱和\ksonline{}三種使用方式，共通的建模與計算步驟統一講解，安裝、檔案管理和遠端執行的差別在相應章節說明。

本書採用中文介面名稱；以英文顯示的欄位保留原名。各計算案例同時給出輸入條件與參考結果，便於操作和比較。

\vspace{\baselineskip}
\noindent{\sffamily\heiti 品牌}

瀚海量子\regmark{}與 KSSOLV Toolbox 為本產品的品牌標識。MATLAB\regmark{} 是 The MathWorks, Inc. 的註冊商標。

\vspace{\baselineskip}
\noindent{\sffamily\heiti 封面}

封面晶體取自隨書矽案例，沿 $[100]$ 方向延展三個常規晶胞。球體表示矽原子，連線連線第一近鄰；青綠色的四條鍵突出金剛石結構中的四面體配位。

\vfill
\noindent{\small © 2026 合肥瀚海量子科技有限公司\par
\noindent Hefei Hanhai Quantum Technology Co., Ltd.}
\clearpage

\noindent{\sffamily\heiti\Large 關於本書}

\vspace{4mm}
\noindent{\sffamily\heiti 適用版本}

本書面向 KSSOLV Toolbox V0.5.0 版本。可在\ui{主頁 → 幫助 → 關於 KSSOLV}中檢視軟體版本。

\vspace{\baselineskip}
\noindent{\sffamily\heiti 建議的閱讀順序}

第一次使用時，可先在第\ref{ch:introduction}章選擇使用方式，再按安裝或線上版章節準備執行環境，隨後閱讀第\ref{ch:workspace}章熟悉工作區。第\ref{ch:first}章的氫分子教程帶讀者完成一次從結構到結果的計算，可按步驟操作並核對每一階段的結果。

準備自己的材料模型時，可按第\ref{ch:modeling}章編輯原子結構，再到第\ref{ch:advanced}章學習材料例項；第\ref{ch:modeling-tools}章提供各項建模工具的查閱說明。計算引數一章解釋常見設定的含義；工作流、體資料與結果章節則用於進一步安排任務和閱讀輸出。不必在第一次操作前讀完所有章節。

隨書的 \file{examples/} 目錄提供結構、專案、復算指令碼與體資料。練習時從指定檔案開始，儲存自己的專案副本；比較數值時，同時核對結構、贗勢和計算條件。

\vspace{\baselineskip}
\noindent{\sffamily\heiti 獲取軟體與幫助}

\begin{itemize}
\item 軟體與發行包：\url{https://github.com/USTC-Hu-Group/KSSOLV-Toolbox}
\item 線上版：\ksonline{} 的瀚海量子\regmark{}官方線上版即將推出，並支援叢集私有化部署。
\item 問題反饋：上述倉庫的 Issues 頁面；也可從應用\ui{幫助}選單進入相關資源。
\end{itemize}

\clearpage
